\documentclass{article}
\textheight = 21cm
\textwidth = 18cm
\topmargin = -2cm
\oddsidemargin= -2cm
\usepackage[T1]{fontenc}
\usepackage[utf8]{inputenc}
\usepackage{graphicx}
\usepackage[spanish]{babel}
\title{Practica01 }
\date {\today }
\author{Diego Mèndez Medina }
\begin{document}   
\begin{center}
  
 
   {\fontsize{23}{0}\selectfont Introducción a las Ciencias de la Computación.}
  \\
  {\fontsize{23}{0}\selectfont Facultad de Ciencias, UNAM. }
  \\
  \includegraphics[width=3cm]{/home/chubak/ICC/ciencias.png}\\
  \\
 {\fontsize{21}{0}\selectfont \fbox { Práctica 01. }} 
 \\
 {\fontsize{20}{0}\selectfont 
  Méndez Medina Diego, Grupo 702.  \\    }
Profesor : 	Salvador López Mendoza. \\	
Ayudante :	Marisol Amezcua. \\
Ayud. Lab. : 	Alejandro Iabin Arteaga Hernández.	
\end{center}
\begin{enumerate}
\item ¿Qué es una clase? \\
Es una plantilla que define la forma de los objetos que ulitizaremos en el programa. \\
\item ¿Cúal es la diferencia entre una clase y un bojeto? \\
Mediante las clases creamos objetos. La clase se codifica y el objeto se crea en el momento que el programa es ejecutado. \\
\item Dar 3 ejemplos de  clases y un par de objetos por cada clase. \\
  \begin{enumerate}
  \item class Estudiante \\
    \begin{itemize}
    \item Diego Méndez Medina  \\
      \item Pedro Pedroza Olvera \\
      \end{itemize}      


     \item class Automovil  \\
    \begin{itemize}
    \item Tesla Model 3 \\
      \item Tesla Model Y \\
    \end{itemize}

     \item class Boligrafo \\
    \begin{itemize}
    \item Bic \\
      \item PaperMate \\
      \end{itemize}  
    \end{enumerate}
  
\item ¿Cuál es el proposito de la documentación? \\
Que cualquiera pueda entender el codigo que escribiste. Tambien aplica para tu yo del futuro.\\
\item ¿Qué ventajas tiene hacer el diseño de un programa antes de codificar? \\
Te permite saber el número de variables y funciones que usara tu programa, o los parametros que necesitaras codificar. Te permite tener la solución o una idea de como hacerlo. \\
\item ¿Qué comando se usa para ejecutar un programa en java ? \\
java NombreDelArchivo.java  \\
\item ¿Qué comando se usa para compilar un programa en java ? \\
javac NombreDelArchivo \\

\end{enumerate}	
\end{document}
